\chapter{练习与反思提示}
\label{ch:exercises}

\section{简答题}

\begin{exercisebox}
\begin{exercise}
请用一段话说明框架惯性与意志能量之间的区别。为什么把两者都压缩成
同一个词“动机”会是一个错误？
\end{exercise}

\begin{exercise}
请给出一个出现在你日常作息中的决策窗口示例。说明为什么这个窗口
持续时间很短。
\end{exercise}

\begin{exercise}
请描述一个目标状态过于模糊、以至于无法进入的案例。障碍分解中的哪一项
捕捉了这种模糊性？
\end{exercise}
\end{exercisebox}

\section{应用建模任务}

\begin{exercisebox}
\begin{exercise}
请对“从被动信息消费转向主动记笔记”的转变进行建模。写出当前状态、
目标状态、主导奖励来源、切换成本，以及一个环境强化因素。
\end{exercise}

\begin{exercise}
请选择一个重复出现的行为失败，并使用第~\ref{ch:interventions}~章中的模板
设计一个一页式干预方案。你的答案必须包含一个决策窗口和一条验证规则。
\end{exercise}
\end{exercisebox}

\section{对比与边界问题}

\begin{exercisebox}
\begin{exercise}
请给出一个案例，其中主要问题并不是当前惯性过高，而是长期目标缺乏清晰度。
说明为什么这个框架仍然有用，但仅凭它本身还不够。
\end{exercise}

\begin{exercise}
在什么情况下，外部承诺会为某个人降低障碍，却为另一个人抬高障碍？
请使用奖励、环境和认知成本的语言来回答。
\end{exercise}
\end{exercisebox}

\section{自我审计工作表}

\begin{templatebox}
\textbf{每周框架惯性审计}
\begin{enumerate}[nosep]
  \item 这周我没能启动的有益转变是什么？
  \item 最终胜出的当前状态是什么？
  \item 哪个障碍项主导了这次失败？
  \item 我错过了哪个决策窗口？
  \item 在下一次尝试之前，我会预先安装哪一个准备性改变？
  \item 下周出现什么证据，才算真实改善？
\end{enumerate}
\end{templatebox}

\section{面向教师的提示使用方式}

本章也可以支持基于讨论的教学：
\begin{itemize}[nosep]
  \item 让学生用同一个模型比较两个失败的转变；
  \item 要求学生举出一个最佳干预主要来自环境设计而非动机提升的案例；
  \item 让学生通过缩小进入动作来重新设计一次失败。
\end{itemize}

\section{推荐阅读}

\begin{readingbox}
\textbf{基础}
\begin{itemize}[nosep]
  \item 在尝试对这些练习做数学扩展之前，先回看记号附录和案例库。
\end{itemize}
\textbf{现代研究}
\begin{itemize}[nosep]
  \item 若后续要把这些练习发展成有研究支撑的问题集，literature map
        是最佳路线。
\end{itemize}
\textbf{前沿}
\begin{itemize}[nosep]
  \item 前沿阅读适合用作研讨课提示，而不适合作为基线评估材料。
\end{itemize}
\end{readingbox}

\begin{summarybox}
练习层把这份手稿从“可阅读的成品”变成“可训练的成品”。这一点很关键，
因为没有练习的教材，无法证明它的概念能够迁移到纸面之外。
\end{summarybox}